\documentclass[11pt]{article}
\usepackage[a4paper,margin=1in]{geometry}
\usepackage{amsmath,amssymb}
\usepackage{booktabs,longtable}
\usepackage{graphicx}
\usepackage[hyphens]{url}
\usepackage[colorlinks=true,linkcolor=black,urlcolor=blue,citecolor=blue]{hyperref}
\setlength{\parskip}{0.55em}\setlength{\parindent}{0pt}
\providecommand{\tightlist}{\setlength{\itemsep}{0.2em}\setlength{\parskip}{0pt}}

\title{Supplement to the Edwards Aquifer J-17 scored test}
\author{\textbf{Amin Abaee}\\ Independent Researcher\\[2pt] ORCID: \href{https://orcid.org/0000-0002-0019-1842}{0000-0002-0019-1842}\\ Email: \href{mailto:amin_abaee@ut.ac.ir}{amin_abaee@ut.ac.ir}}
\date{September 18, 2026}
\begin{document}
\maketitle

\section{Supplement to ``Does a one-pool water-balance model improve
forecasts
of}\label{supplement-to-does-a-one-pool-water-balance-model-improve-forecasts-of}

Edwards Aquifer head? A scored test at J-17''

Section S1 carries the independent replication of the uncertainty layer
of Section 5.3.1; Section S2 carries the truncated-detail climate
fixed-window record of Section 5.4. All values are archived with the
analysis code and re-compute from the registered per-origin forecast
files.

\subsection{S1 Independent replication of the uncertainty
layer}\label{s1-independent-replication-of-the-uncertainty-layer}

An independent replication of this layer is registered with the
repository: the archived script
\texttt{campaign\_e3\_dm\_uncertainty.py}, a second Diebold--Mariano +
moving-block-bootstrap implementation --- HAC lag h − 1 with unweighted
truncation and population-variance scaling, block length h, 20,000
replications, seed 0, deterministic (byte-identical on re-execution) ---
reading the same registered per-origin forecast file. Its verified
output (archived alongside it) reproduces every load-bearing reading:
the M1 retention margin DM z = −0.86 with block CI {[}−1.22, +0.56{]} ft
(bootstrap p = 0.38); the M2m edge over persistence (p = 0.001); the h =
5 climatology win (p = 0.035). It also registers two rows this table
does not carry --- the h = 1 climatology loss (+2.94 ft against
persistence, p = 0.046) and the M2 h = 5 loss (+12.4 ft, p \textless{}
0.001) --- both consistent with the readings above. The two
implementations use different seeds, block lengths, HAC weighting, and
variance conventions, and agree on every conclusion: the 0.39-ft AR(1)
retention margin is statistically indistinguishable from zero, the M2m
edge is the only one-year margin that separates from noise, and the
five-year climatology win survives.

\subsection{S2 Climate-informed modules on the remaining fixed
windows}\label{s2-climate-informed-modules-on-the-remaining-fixed-windows}

The remaining fixed windows complete the same record: drawdown 1951--56,
M2\_Renso 24.30, M2\_Rprecip 28.67, M2\_combo 24.70, precipitation
oracle 18.16 ft; pre-permit wet 1991--95, M2\_Rprecip 22.03, M2\_Renso
23.03, M2\_Rar 23.94, M2\_combo 25.71, oracle 10.98 ft; critical-period
era 2015--23, M2\_Rprecip 14.52, M2\_Rar 14.67, M2\_Renso 16.01,
M2\_combo 16.57, oracle 9.75 ft --- on this one window the two
climate-informed modules M2\_Rprecip and M2\_Rar edge past M1 by about
one foot (14.52 and 14.67 versus 15.62 ft). On the recharge target
itself the fixed-window scores are an order of magnitude coarser on
every window (climate modules 199--937; precipitation oracle 80.7--487,
against the 556 × 10³ acre-ft climatology scale of the rolling record),
so the marginal head advantage is a window-specific result, not a
recharge forecast, and it does not enter the abstract.

\end{document}
